\documentclass[12pt,a4paper,oneside]{report}

% ─── Packages ─────────────────────────────────────────────────────────────────
\usepackage[utf8]{inputenc}
\usepackage[T1]{fontenc}

\usepackage{amsmath,amssymb,amsfonts,amsthm}
\usepackage{graphicx}
\usepackage{geometry}
\usepackage{setspace}
\usepackage{fancyhdr}
\usepackage{titlesec}
\usepackage{tocloft}
\usepackage{booktabs}
\usepackage{multirow}
\usepackage{array}
\usepackage{longtable}
\usepackage{tabularx}
\usepackage{url}
\usepackage{xcolor}
\usepackage{caption}
\usepackage{tikz}
\usepackage{pgfplots}
\pgfplotsset{compat=1.18}
\usetikzlibrary{shapes.geometric,arrows.meta,positioning,fit,backgrounds,calc,shapes.multipart,chains}
\usepackage{enumitem}
\usepackage{float}
\usepackage{hyperref}

% ─── Page geometry ────────────────────────────────────────────────────────────
\geometry{a4paper,top=1in,bottom=1in,left=1.5in,right=1in,headheight=15pt}

% ─── Spacing ──────────────────────────────────────────────────────────────────
\onehalfspacing

% ─── Hyperref ─────────────────────────────────────────────────────────────────
\hypersetup{colorlinks=true,linkcolor=black,citecolor=black,
  urlcolor=blue!60!black,
  pdftitle={ZK-Auth Thesis Synopsis},pdfauthor={Abhay Dandge}}

% ─── Header/Footer ────────────────────────────────────────────────────────────
\pagestyle{fancy}
\fancyhf{}
\fancyfoot[C]{\small\thepage}
\renewcommand{\headrulewidth}{0pt}
\fancypagestyle{plain}{
  \fancyhf{}
  \fancyfoot[C]{\small\thepage}
  \renewcommand{\headrulewidth}{0pt}
}

% ─── TOC Dots for Chapters ────────────────────────────────────────────────────
\renewcommand{\cftchapleader}{\cftdotfill{\cftdotsep}}

% ─── Initial Chapter/section format (Without lines for Front Matter) ──────────
\titleformat{\chapter}[display]
  {\normalfont\Large\bfseries\raggedleft}{CHAPTER~\thechapter}{12pt}{\large\MakeUppercase}
\titlespacing*{\chapter}{0pt}{-10pt}{30pt}

\titleformat{name=\chapter,numberless}[block]
  {\normalfont\Large\bfseries\raggedleft}{}{0pt}{\large\MakeUppercase}

\titleformat{\section}{\normalfont\large\bfseries}{\thesection}{1em}{}
\titleformat{\subsection}{\normalfont\normalsize\bfseries}{\thesubsection}{1em}{}

% ─── TikZ styles ──────────────────────────────────────────────────────────────
\tikzset{
  block/.style={rectangle,draw=black!70,fill=blue!8,rounded corners=3pt,
    minimum width=2.6cm,minimum height=0.8cm,text width=2.4cm,align=center,
    font=\small},
  actor/.style={rectangle,draw=black!80,fill=gray!10,rounded corners=4pt,
    minimum width=2.8cm,minimum height=1.0cm,text width=2.6cm,align=center,
    font=\small\bfseries},
  storage/.style={cylinder,draw=black!70,fill=yellow!10,
    shape border rotate=90,minimum width=1.6cm,
    minimum height=0.7cm,text centered,font=\scriptsize},
  arrow/.style={-{Stealth[length=4pt]},draw=black!70,line width=0.6pt},
  label/.style={font=\scriptsize, align=center}
}

% ─── Theorem environments ─────────────────────────────────────────────────────
\newtheorem{definition}{Definition}[chapter]
\newtheorem{theorem}{Theorem}[chapter]

%==============================================================================
\begin{document}
%==============================================================================

%------------------------------------------------------------------------------
% TITLE PAGE
%------------------------------------------------------------------------------
\begin{titlepage}
\vspace*{\fill}
\begin{center}
{\Large\bfseries
ZERO-KNOWLEDGE PROOF-BASED AUTHENTICATION FRAMEWORK\\[4pt]
WITH SELECTIVE DISCLOSURE AND BEHAVIORAL BIOMETRICS\\[4pt]
FOR CLOUD COMPUTING
}\\[0.6cm]
\vfill
{\large A}\\[0.2cm]
{\Large\bfseries SYNOPSIS}\\[0.3cm]
{\normalsize\itshape
Submitted in partial fulfillment of the requirements for the award of degree of}\\[0.3cm]
{\large\bfseries MASTERS OF TECHNOLOGY}\\[0.2cm]
{\normalsize\itshape In}\\[0.2cm]
{\large\bfseries ARTIFICIAL INTELLIGENCE}\\[0.4cm]
\vfill
{\normalsize\itshape Submitted by}\\[0.2cm]
{\large\bfseries ABHAY DANDGE}\\[0.2cm]
{\normalsize Scholar Number: \textbf{24215011117}}\\[0.4cm]
\vfill
{\normalsize\itshape Under the Guidance of}\\[0.2cm]
{\normalsize\bfseries Dr.~Namita Tiwari}\\
{\normalsize Assistant Professor, Department of CSE, MANIT Bhopal}\\[0.2cm]
{\normalsize\bfseries Dr.~Meenu Chawla}\\
{\normalsize Professor, Department of CSE, MANIT Bhopal}\\[0.5cm]
\vfill
\includegraphics[width=2.5cm]{manit_logo.png}\\[0.2cm]
{\normalsize\bfseries Centre of Artificial Intelligence}\\[0.2cm]
{\large\bfseries MAULANA AZAD NATIONAL INSTITUTE OF TECHNOLOGY BHOPAL}\\[0.1cm]
{\normalsize\bfseries DECEMBER 2025}
\end{center}
\vspace*{\fill}
\end{titlepage}

%------------------------------------------------------------------------------
% FRONT MATTER
%------------------------------------------------------------------------------
\clearpage
\pagenumbering{roman}
\setcounter{page}{2}

% DECLARATION
\clearpage
\phantomsection
\addcontentsline{toc}{chapter}{Declaration}
\begin{center}
{\large\bfseries MAULANA AZAD NATIONAL INSTITUTE OF TECHNOLOGY BHOPAL}\\[0.1cm]
{\normalsize (An Institute of National Importance)}\\[0.1cm]
{\normalsize Bhopal - 462003}\\[0.4cm]
\includegraphics[width=2.5cm]{manit_logo.png}\\[0.4cm]
{\normalsize\bfseries Centre of Artificial Intelligence}\\[0.6cm]
{\Large\bfseries\underline{Declaration}}
\end{center}

\vspace{0.5cm}
\noindent I hereby declare that the work presented in this synopsis, summarizing
the dissertation entitled \textbf{``Zero-Knowledge Proof-Based Authentication
Framework with Selective Disclosure and Behavioral Biometrics for Cloud
Computing''}, is submitted in partial fulfillment of the requirements for the
award of the degree of \textbf{Masters of Technology} in \textbf{Artificial
Intelligence}. It is an authentic summary of my own original work carried out
from July~2024 to December~2025 under the guidance of \textbf{Dr.~Namita
Tiwari}, Assistant Professor, and co-supervisor \textbf{Dr.~Meenu Chawla},
Professor, Department of CSE, MANIT Bhopal.

Wherever contributions of others are involved, every effort is made to indicate
this clearly with due reference to the literature. The matter embodied in this synopsis has not been presented or submitted for any other degree.

\vspace{2.5cm}
\begin{flushright}
\rule{6cm}{0.4pt}\\[5pt]
\textbf{Abhay Dandge}\\
Scholar Number: 24215011117\\
M.Tech (Artificial Intelligence)\\
MANIT, Bhopal
\end{flushright}

% CERTIFICATE
\clearpage
\phantomsection
\addcontentsline{toc}{chapter}{Certificate}
\begin{center}
{\large\bfseries MAULANA AZAD NATIONAL INSTITUTE OF TECHNOLOGY BHOPAL}\\[0.1cm]
{\normalsize (An Institute of National Importance)}\\[0.1cm]
{\normalsize Bhopal - 462003}\\[0.4cm]
\includegraphics[width=2.5cm]{manit_logo.png}\\[0.4cm]
{\normalsize\bfseries Centre of Artificial Intelligence}\\[0.6cm]
{\Large\bfseries\underline{Certificate}}
\end{center}

\vspace{0.5cm}
\noindent This is to certify that the synopsis entitled \textbf{``Zero-Knowledge
Proof-Based Authentication Framework with Selective Disclosure and Behavioral
Biometrics for Cloud Computing''} submitted by \textbf{Abhay Dandge},
Scholar No.~24215011117, in partial fulfillment of the requirements for the
degree of \textbf{Masters of Technology} in Artificial Intelligence, summarizes
authentic work carried out under our supervision and guidance from
July~2024 to December~2025.

\vspace{3cm}
\noindent
\begin{minipage}[t]{0.45\textwidth}
\rule{5.5cm}{0.4pt}\\[3pt]
\textbf{Dr.~Namita Tiwari}\\
Supervisor\\
Assistant Professor, Dept. of CSE\\
MANIT, Bhopal
\end{minipage}
\hfill
\begin{minipage}[t]{0.45\textwidth}
\raggedleft
\rule{5.5cm}{0.4pt}\\[3pt]
\textbf{Dr.~Meenu Chawla}\\
Co-supervisor\\
Professor, Dept. of CSE\\
MANIT, Bhopal
\end{minipage}

% ─── Redefine Chapter Formats for Line Under Headings ─────────────────────────
\titleformat{name=\chapter,numberless}[block]
  {\normalfont\Large\bfseries\raggedleft}{}{0pt}{\large\MakeUppercase}[\vspace{1ex}\titlerule]

\titleformat{\chapter}[display]
  {\normalfont\Large\bfseries\raggedleft}{CHAPTER~\thechapter}{12pt}{\large\MakeUppercase}[\vspace{1ex}\titlerule]
% ──────────────────────────────────────────────────────────────────────────────

% ABSTRACT (ENGLISH ONLY, ONE PAGE)
\clearpage
\chapter*{Abstract}
\addcontentsline{toc}{chapter}{Abstract}

Cloud computing has fundamentally altered digital service delivery, yet
authentication security remains a critical vulnerability as infrastructure
scales. Traditional mechanisms---passwords, multi-factor authentication, and
public key infrastructure---rely on a \textit{verify-by-reveal} paradigm in
which sensitive credentials must be transmitted or stored by the
authenticating party, creating attack surfaces repeatedly exploited through
database breaches, credential stuffing, and man-in-the-middle attacks.
Delegation alternatives such as OAuth~2.0 and SAML do not eliminate this
flaw; they merely centralize it into large-scale identity providers.

This thesis introduces \textsc{ZK-Auth}, a cryptographic infrastructure
resolving the verify-by-reveal problem. \textsc{ZK-Auth} is best described
as an \textbf{OAuth 2.0 Authorization Server in which the standard
``user authenticates'' step is implemented via a zero-knowledge proof
rather than a password} --- the same architectural pattern used by
``Sign in with Apple,'' where OAuth/OIDC governs relying-party delegation
(which application is asking, for what scope, where the user is redirected
afterward) while a non-password authenticator (Face~ID there; a Groth16
proof here) governs the identity check underneath. This thesis therefore
does not claim that ZKP alone replaces every authentication primitive;
rather, ZKP replaces the password/credential layer that OAuth has always
treated as a delegated, swappable component.

Three mechanisms realize this design. \textit{First}, the OAuth
Authorization Code grant's authentication step is implemented using
Groth16 zero-knowledge succinct non-interactive arguments of knowledge
(zk-SNARKs) over the BN254 elliptic curve, compiled via Circom~2.x and
deployed as browser-executable WebAssembly, allowing users to prove
knowledge of a registration secret without ever transmitting it; the
resulting proof authorizes a short-lived OAuth authorization code rather
than a password-backed session token. Poseidon hash-based nullifiers
prevent replay attacks while keeping server-side verification to tens of
milliseconds. \textit{Second}, selective credential disclosure
is achieved via depth-8 Poseidon Merkle trees integrated with W3C
Verifiable Credentials, enabling holders to prove individual attribute
predicates (e.g., age $\geq$ 18) without exposing raw attribute values or
salts. \textit{Third}, recognizing that point-in-time authentication
cannot protect against session hijacking, \textsc{ZK-Auth} incorporates a
two-layer LSTM behavioral biometric model providing continuous,
background session-risk scoring from telemetry, dynamically triggering
cryptographic step-up re-authentication on anomalous behavior. The
architecture further extends to a multi-tenant institutional issuance
platform in which registered institutions receive unique Ed25519 signing
keypairs, analogous to a distributed PKI certificate-authority hierarchy.

Synthesizing and extending two peer-reviewed M.Tech
publications~\cite{dandge2025zkpcloud,prasad2025abereview}, \textsc{ZK-Auth}
is implemented end-to-end --- including a working OAuth~2.0 Authorization
Code flow with PKCE support --- and empirically evaluated across four
heterogeneous deployment environments: a Node.js server, a Chrome browser
client, and two Android devices of differing hardware tier. Across 120
consecutive server-side trials, the system achieves a median end-to-end
authentication latency of 132.5\,ms---23\% faster than the Argon2id
password-hashing step alone in a comparable conventional system---with a
compact $\sim$723-byte proof; client-side proof generation across all
measured environments remains within sub-second bounds, from 61.9\,ms
(Node.js) to 361.0\,ms (mid-tier Android tablet). A dedicated baseline
comparison against seven existing authentication paradigms (Password,
MFA/TOTP, PKI smart-card, pure OAuth~2.0/OIDC, SAML, WebAuthn/FIDO2, and
behavioral-biometric-only systems) is presented with both qualitative
feature analysis and literature-reported quantitative figures. The
behavioral biometric subsystem, trained and evaluated on the public
Balabit Mouse Dynamics Challenge dataset rather than synthetic telemetry,
achieves an AUC-ROC of 0.816 with a 6.15\% false-acceptance rate. A formal
extraction-based reduction further establishes that the circuit's
shared-witness construction preserves Groth16's soundness guarantee for
the specific Poseidon-bound commitment and nullifier bindings used in this
system. By completely eliminating server-side storage of personally
identifiable information, \textsc{ZK-Auth} achieves structural compliance
with India's Digital Personal Data Protection Act 2023 and the EU GDPR.

% TABLE OF CONTENTS
%------------------------------------------------------------------------------
\clearpage
\tableofcontents

\clearpage
\listoftables
\addcontentsline{toc}{chapter}{List of Tables}

\clearpage
\listoffigures
\addcontentsline{toc}{chapter}{List of Figures}

%==============================================================================
% MAIN MATTER
%==============================================================================
\clearpage
\pagenumbering{arabic}
\setcounter{page}{1}

%==============================================================================
\chapter{Introduction}
%==============================================================================

\section{Background and Motivation}

Cloud computing has transformed digital service delivery, but authenticating
users in environments where infrastructure is not under direct organizational
control remains a critical security challenge. Traditional authentication
mechanisms were designed for environments where the authenticating server is
trusted; in cloud deployments this assumption often does not hold. The
\textit{verify-by-reveal} paradigm requires credentials to be transmitted to
or stored by the authenticating party, creating attack surfaces that
cloud-scale breaches repeatedly exploit.

Zero-Knowledge Proofs (ZKPs) offer a mathematically rigorous resolution: a
\textit{prover} convinces a \textit{verifier} that a statement is true
without revealing any information beyond its truth~\cite{gmr89}. Applied to
authentication, a user proves knowledge of a registration secret without
transmitting it---eliminating the primary attack surface of credential
theft. Despite this theoretical appeal, ZKPs have seen limited production
adoption due to four challenges this thesis addresses: (1)~proof-generation
complexity and the lack of cross-environment latency characterization in
prior work; (2)~the absence of attribute-level selective disclosure in
basic ZKP authentication; (3)~the inability of one-time proofs to prevent
post-authentication session hijacking; and (4)~a frequent lack of empirical
benchmarking or constraint-level security reduction in published ZKP
authentication systems.

A fifth, more architectural observation also motivates this thesis: most
ZKP-authentication proposals in the literature implicitly assume ZKP
\textit{replaces} the relying-party delegation layer that systems like
OAuth~2.0 already provide, rather than \textit{slotting into} it as the
authentication primitive underneath. \textsc{ZK-Auth} is designed around
the latter view: it is implemented as an OAuth~2.0 Authorization Server
(Chapter~3) in which the conventional password-based login step of the
Authorization Code grant is replaced by a Groth16 proof, so institutions
can adopt ZKP-based authentication without discarding existing
OAuth-based relying-party integrations.

\section{Published Research Underpinning This Thesis}

This thesis synthesizes and substantially extends two peer-reviewed
publications: \textbf{Dandge, A., Prasad, S., Tiwari, N., and Chawla, M.},
``Zero-Knowledge Proof for Authentication in Cloud Computing: A
Review''~\cite{dandge2025zkpcloud} (Springer Nature, FrontSci~2025), which
established the architectural taxonomy of ZKP cloud-authentication
deployment patterns and motivated the Groth16 design choice used
throughout; and \textbf{Prasad, S., Dandge, A., Tiwari, N., and Chawla,
M.}, ``A Comprehensive Review of Zero-Knowledge Proofs in Attribute-Based
Encryption Frameworks''~\cite{prasad2025abereview} (accepted, ICNDIA~2026,
IEEE), which informed the selective-disclosure design developed in
Chapter~3. Both publications surveyed the literature; this thesis is the
first work in this research programme to implement, deploy, and
empirically benchmark a complete system.

\section{Research Contributions}

\begin{enumerate}
  \item A complete implementation of Groth16 zk-SNARK authentication over
        BN254 using Circom-compiled circuits, deployed as browser-
        executable WebAssembly with Node.js server-side verification,
        wrapped in a working \textbf{OAuth~2.0 Authorization Code flow
        with PKCE} so the proof substitutes for the conventional
        password-login step rather than replacing relying-party
        delegation entirely.
  \item A depth-8 Poseidon Merkle tree commitment scheme for W3C VC
        selective attribute disclosure supporting GTE/LTE/EQ predicate
        proofs without attribute value exposure, with an ablation analysis
        justifying the depth-8 design point.
  \item A two-layer LSTM behavioral biometric model on a 50-event sliding
        window, trained and evaluated on the public Balabit Mouse
        Dynamics Challenge dataset, with automated ZKP step-up
        re-authentication.
  \item A multi-tenant institutional credential issuance platform with
        per-institution Ed25519 keypair generation and W3C DID
        publication.
  \item Cross-environment empirical validation across four deployment
        environments, together with a formal extraction-based soundness
        reduction for the implemented circuit, and a baseline comparison
        against seven existing authentication paradigms.
\end{enumerate}

\section{Organisation of the Synopsis}

Chapter~2 summarizes the literature and theoretical foundations. Chapter~3
presents the system architecture, ZKP protocol, and advanced subsystems.
Chapter~4 summarizes the implementation and key experimental results.
Chapter~5 concludes with the principal findings and future scope.

%==============================================================================
\chapter{Literature Survey and Theoretical Background}
%==============================================================================

\section{Zero-Knowledge Proof Systems}

Zero-Knowledge Proofs were introduced by Goldwasser, Micali, and
Rackoff~\cite{goldreich94}, establishing completeness, soundness, and
zero-knowledge as foundational properties. Non-interactive variants via a
common reference string~\cite{bfm88} and the Fiat--Shamir
heuristic~\cite{fiatshamir} enabled practical succinct non-interactive
arguments of knowledge (SNARKs), culminating in Pinocchio~\cite{parno13}
and the quadratic arithmetic program formulation~\cite{gennaro13}.
Groth16~\cite{groth16}, the construction used throughout this thesis,
achieves constant-size proofs (three group elements) with
$\mathcal{O}(1)$ pairing-based verification---the most efficient known
SNARK under the generic group model for a fixed circuit---and has been
deployed at scale in Zcash~\cite{zcash14}. Alternatives such as
PLONK~\cite{plonk} (universal setup) and STARKs~\cite{starks} (transparent
setup) trade Groth16's compact proof size for setup flexibility, at a
proof-size cost incompatible with this thesis's latency- and bandwidth-
sensitive authentication context.

\section{ZKP for Cloud Authentication}

Our prior review~\cite{dandge2025zkpcloud} surveyed 28 papers (2019--2025),
distinguishing stand-alone ZKP deployment (the cloud service itself acts
as verifier)~\cite{khandait2024,pathak2021,butt2024} from ledger-anchored
deployment (smart-contract verification)~\cite{hou2022,verma2024,naidu2023}.
Table~\ref{tab:zkp_families} reproduces this review's central comparative
finding, establishing Groth16 as optimal for high-traffic, latency-
sensitive cloud APIs---the basis for this thesis's design choice.

\begin{table}[H]
\renewcommand{\arraystretch}{1.2}
\caption{ZKP Family Comparison for Cloud Authentication~\cite{dandge2025zkpcloud}}
\label{tab:zkp_families}
\centering
\scriptsize
\begin{tabular}{lccc}
\toprule
\textbf{ZKP Family} & \textbf{Setup} & \textbf{Proof Size} & \textbf{Verif. Cost} \\
\midrule
$\Sigma$-protocols (Schnorr) & None    & Small         & Low \\
zk-SNARKs (Groth16)          & Trusted & $\approx$288B & $\mathcal{O}(1)$ \\
Universal SNARKs (PLONK)     & Universal & Small       & $\mathcal{O}(1)$ \\
zk-STARKs                    & Transparent & $>$40KB   & Low \\
\midrule
\textbf{ZK-Auth uses:} & \multicolumn{3}{c}{\textbf{Groth16}} \\
\bottomrule
\end{tabular}
\end{table}

\section{Selective Disclosure and Behavioral Biometrics}

BBS+ signatures~\cite{bbs} and SD-JWT~\cite{sdjwt} support attribute-level
disclosure but expose the revealed attribute value in cleartext to the
verifier; CL credentials~\cite{cl} provide genuine zero-knowledge
disclosure but predate SNARK-friendly hash functions, yielding larger
proofs than a Poseidon/Groth16 combination. The W3C VC Data
Model~\cite{w3cvc} and DID specification~\cite{w3cdid} provide the
disclosure-mechanism-agnostic standards framework adopted in this thesis.
For continuous authentication, behavioral biometrics via keystroke
dynamics~\cite{shepherd95}, mouse patterns~\cite{nakkabi11}, and deep
sequence models such as TypeNet~\cite{typenet} demonstrate that recurrent
architectures extract stable behavioral signatures from interaction
telemetry; commercial systems such as TypingDNA~\cite{typingdna} confirm
production viability. These systems predominantly use behavioral
biometrics as a \textit{primary} authenticator, whereas \textsc{ZK-Auth}
positions the LSTM layer as a \textit{secondary}, post-ZKP
session-continuity check.

\section{Cryptographic Primitives}

Poseidon~\cite{poseidon} is a zk-SNARK-optimized hash achieving
$\approx$128-bit collision resistance with substantially fewer R1CS
constraints per invocation than a bit-oriented hash, directly determining
the compact circuit size measured in Chapter~4. Ed25519~\cite{bernstein2011}
provides 128-bit-secure digital signatures used at the institutional-
issuer layer. LSTM networks~\cite{hochreiter97} maintain gated hidden and
cell states that mitigate the vanishing-gradient problem limiting plain
recurrent networks on long sequences, motivating their use for the
50-event behavioral window developed in Chapter~3.

%==============================================================================
\chapter{System Architecture and Subsystems}
%==============================================================================

\section{Threat Model and Trust Architecture}

The adversary model permits passive network observation, active
modification of server-side records excluding protected key material, and
full knowledge of circuit structure, while excluding cryptographic breaks
of BN254, Poseidon, or Ed25519. Security goals are: authentication
soundness, zero-knowledge, nullifier replay prevention, selective
disclosure privacy, and bounded-window session-continuity assurance.

\textsc{ZK-Auth} implements the W3C VC trust triangle
(Fig.~\ref{fig:trust_model}) with the Platform as root authority. The
Platform authorizes institutions without issuing credentials directly;
Issuers store only Poseidon Merkle roots; Holders generate Groth16 proofs
locally; Verifiers receive only boolean predicate results.

\begin{figure}[H]
\centering
\resizebox{0.55\columnwidth}{!}{%
\begin{tikzpicture}[node distance=1.6cm and 1.6cm]
  % Nodes
  \node[actor, fill=orange!12] (platform) {ZK-Auth Platform};
  \node[actor, fill=blue!10, below left=of platform] (issuer) {Issuer};
  \node[actor, fill=purple!10, below right=of platform] (verifier) {Verifier};
  \node[actor, fill=green!10, below right=of issuer] (holder) {Holder};

  % Edges (Straight arrows for simple connections)
  \draw[arrow] (platform) -- node[label, above left] {Ed25519 key} (issuer);
  \draw[arrow] (issuer) -- node[label, below left] {Issue VC} (holder);
  \draw[arrow] (holder) -- node[label, below right] {ZKP proof} (verifier);
  
  % Curved lines to prevent overlapping arrows and text
  \draw[arrow, dashed, bend left=12] (verifier) to node[label, left, align=right, xshift=-0.1cm] {DID\\Resolve} (platform);
  \draw[arrow, bend left=12] (platform) to node[label, right, xshift=0.1cm] {Public Key} (verifier);

\end{tikzpicture}%
}
\caption{Three-actor trust model. The ZK-Auth Platform acts as the root authority, authorizing institutions (Issuers) with unique keypairs. Holders receive credentials and submit proofs to Verifiers.}
\label{fig:trust_model}
\end{figure}

\section{Circuit Design and Authentication Protocol}

The authentication circuit \texttt{auth.circom} compiles to a Rank-1
Constraint System (R1CS) with \textbf{932 constraints and 935 wires}
(\texttt{snarkjs r1cs info}). With private input $s \in \mathbb{Z}_r$
(registration secret) and public input $n \in \mathbb{Z}_r$ (challenge
nonce), the circuit produces:
\begin{align}
  c &:= \text{Poseidon}(s) &&\text{(commitment root)} \\
  \eta &:= \text{Poseidon}(s \| n) &&\text{(nullifier hash)}
\end{align}
The proving system is Groth16 over BN254, with trusted setup via the
Hermez Network's public Powers of Tau ceremony ($2^{15}$ constraint
ceiling, leaving $31{,}836$ constraints of headroom). At registration, the
client computes $c$ locally and sends only $c$; the secret $s$ never
leaves the client. At authentication, the client requests a challenge
nonce, computes a Groth16 proof via \texttt{groth16.fullProve}, and
submits it; the server verifies the proof, confirms the commitment, and
atomically checks/inserts the nullifier in Redis to prevent replay before
issuing session tokens. Every verification call is padded to a minimum of
50\,ms to eliminate timing side-channels.

\begin{figure}[H]
\centering
\resizebox{0.65\columnwidth}{!}{%
\begin{tikzpicture}[node distance=0.8cm and 0.4cm]
  % Client
  \node[block, fill=blue!10, minimum width=6cm] (clients) {Clients (Web / Mobile)};
  
  % API/Gateway layer
  \node[block, fill=green!8, below left=0.8cm and -1.2cm of clients] (api) {Node.js API};
  \node[block, fill=green!8, below right=0.8cm and -1.2cm of clients] (ws) {WebSockets};
  
  % Compute Layer
  \node[block, fill=purple!8, below=of api] (ml) {LSTM ML (gRPC)};
  \node[block, fill=orange!8, below=of ws] (zkp) {ZKP Verify};
  
  % Storage Layer
  \node[storage, fill=yellow!12, below=0.8cm of ml] (pg) {PostgreSQL};
  \node[storage, fill=yellow!12, right=0.4cm of pg] (ts) {TimescaleDB};
  \node[storage, fill=red!10, below=0.8cm of zkp] (redis) {Redis};
  
  % Edges
  \draw[arrow] (clients.195) -- (api.north);
  \draw[arrow, dashed] (clients.345) -- (ws.north);
  
  % API down to Compute/Storage (using shifted anchors on API's bottom edge to prevent crossing)
  \draw[arrow] (api.south) -- (ml.north);
  \draw[arrow] ([xshift=0.4cm]api.south) -- (zkp.north west);
  \draw[arrow] ([xshift=0.8cm]api.south) -- (redis.north west);
  
  % Route API to PG by going out the left side and down (bypasses ML block)
  \draw[arrow] (api.west) -- ++(-0.5,0) |- (pg.west);
  
  % ML to TS
  \draw[arrow] (ml.south east) -- (ts.north);
  
  % Route WS to Redis by going out the right side and down (bypasses ZKP block)
  \draw[arrow] (ws.east) -- ++(0.5,0) |- (redis.east);

\end{tikzpicture}%
}
\caption{ZK-Auth system component architecture. Clients interface with the Node.js API (REST) and WebSockets. The compute tier handles ML risk scoring and ZKP verification. Data is persisted in PostgreSQL, while Redis manages ephemeral session nonces.}
\label{fig:system_overview}
\end{figure}

\section{ZK-Auth as an OAuth 2.0 Authorization Server}
\label{sec:oauth_zkp_synopsis}

A precise architectural statement matters here: \textsc{ZK-Auth} is not
``pure ZKP authentication.'' It is structurally an \textbf{OAuth~2.0
Authorization Server implementing the RFC~6749 Authorization Code
grant}, in which the conventional password-based login step of that
grant is replaced by the Groth16 challenge/verify exchange above. OAuth
governs \textit{which} relying party is asking, \textit{what} scope it
requests, and \textit{where} the user is redirected; the ZKP layer
governs \textit{how} the user proves identity within that delegation
contract --- the same pattern used by ``Sign in with Apple,'' which
pairs OAuth/OIDC with Face~ID/Touch~ID instead of a password.

Concretely, the relying party registers once as an OAuth client
(\texttt{client\_id} + hashed \texttt{client\_secret}), redirects the
user to \texttt{GET /oauth/authorize}, and the frontend renders the
same ZKP login widget described above rather than a password form. On
successful, non-replayed proof verification, the server mints a
short-lived (60-second), single-use authorization code bound to the
specific nullifier hash that authorized it --- not a session token
directly. The relying party then exchanges this code at
\texttt{POST /oauth/token} (with PKCE \texttt{code\_verifier} support
for public clients) for a real access/refresh token pair, issued via
the same session-issuance path used by direct, non-OAuth logins.
Authorization-code reuse --- a signal of interception --- triggers the
same defense-in-depth response as a detected session hijack: all
sessions for the affected user are revoked.

\section{Selective Disclosure and Behavioral Continuity}

For attribute $a_i$ with value $v_i$ and random salt $\sigma_i$, the leaf
commitment is $\ell_i := \text{Poseidon}(\text{enc}(v_i) \| \sigma_i)$,
aggregated into a depth-8 Poseidon Merkle root $\rho$ (256-leaf capacity).
The disclosure circuit proves Merkle-path validity, attribute-commitment
consistency, and a GTE/LTE/EQ predicate simultaneously; the verifier
receives only $(\rho, \text{predicate}, \tau, \text{result})$, never the
underlying value or salt.

For session continuity, six telemetry features (mouse velocity, key
dwell, scroll delta, touch pressure, event type, inter-event gap) feed a
two-layer LSTM ($\text{LSTM}_{128} \to \text{LSTM}_{64} \to$ Dense~32
$\to$ sigmoid) over a 50-event sliding window, with EMA-smoothed risk
scoring. At risk $r \geq 0.90$, the server forces a fresh ZKP challenge
within a bounded window, closing the session-hijacking gap that one-time
authentication leaves open.

\section{Multi-Tenant Institutional Issuance}

The platform implements a PKI-analogous root-of-trust: each registered
institution receives a unique Ed25519 keypair and a DID resolving to
\texttt{did:web:\{slug\}.zk-auth.io}, then signs W3C VCs embedding the
holder's Merkle root. The private key is delivered once at onboarding and
never persisted by the platform, so a platform-side breach cannot be used
to forge institutional credentials.

%==============================================================================
\chapter{Implementation and Key Results}
%==============================================================================

\section{Experimental Setup}

Server and browser experiments use an Apple MacBook Air M4 running the
full stack (Node.js/Express, PostgreSQL, TimescaleDB, Redis) locally, with
\texttt{snarkjs} v0.7.4 for proof generation/verification. Mobile
measurements use two physical Android devices on the same network: a
Samsung Galaxy~S23 (flagship SoC) and a Samsung Galaxy Tab~S9~FE+
(mid-tier SoC), both running the production Flutter client with the
circuit executed inside a hidden WebView.

\section{Cross-Environment ZKP Performance}

Table~\ref{tab:zkp_results} reports proof-generation and verification
latency across all four measured environments ($n=120$ Node.js trials,
$n=30$ browser trials, $n=15$ per mobile device). Fig.~\ref{fig:env_comparison}
visualizes the comparison directly.

\begin{table}[H]
\renewcommand{\arraystretch}{1.2}
\caption{ZKP Performance Across Environments (median values; server
verify includes the 50\,ms timing pad).}
\label{tab:zkp_results}
\centering
\small
\begin{tabular}{lc}
\toprule
\textbf{Environment} & \textbf{Median proof gen. (ms)} \\
\midrule
Node.js (server)              & 61.93 \\
Chrome browser (Web Worker)   & 153.4 \\
Galaxy S23 (Android WebView)  & 254.0 \\
Tab S9 FE+ (Android WebView)  & 361.0 \\
\midrule
\textbf{Total E2E login (Node.js)} & \textbf{132.51} \\
\bottomrule
\end{tabular}
\end{table}

\begin{figure}[H]
\centering
\begin{tikzpicture}
\begin{axis}[
  width=9cm, height=5.4cm,
  ybar, bar width=16pt,
  symbolic x coords={Node.js, Browser, Galaxy S23, Tab S9 FE+},
  xtick=data, x tick label style={font=\small, rotate=15, anchor=east},
  ylabel={Proof generation (ms)}, label style={font=\small},
  ymin=0, ymax=420,
  ymajorgrids, grid style={line width=0.1pt, draw=gray!25},
  nodes near coords, every node near coord/.style={font=\small},
  axis line style={line width=0.4pt}, tick label style={font=\small},
]
\addplot[fill=blue!55, draw=blue!70!black] coordinates {
  (Node.js,61.93) (Browser,153.4) (Galaxy S23,254.0) (Tab S9 FE+,361.0)
};
\end{axis}
\end{tikzpicture}
\caption{Median proof-generation time across all four measured
environments for the same auth circuit (932 constraints).}
\label{fig:env_comparison}
\end{figure}

The full end-to-end login (132.51\,ms median) is \textbf{23\% faster than
Argon2id} hashing alone (172.82\,ms median, $n=30$), the dominant
cryptographic cost of a conventional password system---demonstrating that
ZKP authentication is not a latency trade-off but a measured improvement.

\section{Behavioral Biometric Performance}

The LSTM model, trained on the public Balabit Mouse Dynamics Challenge
dataset (108{,}111 training windows, 27{,}027 held-out test windows; 50-event
sliding window, stride~25), achieves the results in
Table~\ref{tab:lstm_results}.

\begin{table}[H]
\renewcommand{\arraystretch}{1.2}
\caption{LSTM Behavioral Model Performance (threshold = 0.75).}
\label{tab:lstm_results}
\centering
\begin{tabular}{lc}
\toprule
\textbf{Metric} & \textbf{Value} \\
\midrule
AUC-ROC                  & 0.8157 \\
False Acceptance Rate    & 6.15\% \\
False Rejection Rate     & 56.68\% \\
F1-score (weighted avg.) & 0.8711 \\
Inference time           & 15.83\,ms \\
\bottomrule
\end{tabular}
\end{table}

The conservative 0.75 threshold minimizes false acceptance---the
security-critical error---at the cost of elevated false rejection, which
only adds re-authentication friction rather than a security exposure.

\section{Ablation Summary}

A depth-8 Merkle tree (256-leaf capacity) was selected as the smallest
depth comfortably exceeding realistic per-credential attribute counts
while keeping disclosure proving time near 173\,ms (measured); doubling
to depth-16 roughly doubles proving cost for no benefit at the
per-individual-credential granularity this system targets. Similarly, the
50-event LSTM window balances detection-lag against signal stability,
achieving AUC-ROC~0.816 at 15.83\,ms inference---negligible relative to
the 132.5\,ms end-to-end authentication latency it complements.

\section{System Comparison}

\begin{table}[H]
\renewcommand{\arraystretch}{1.2}
\caption{Comparison with Existing Authentication Systems.
\checkmark~=supported; $\approx$~=partial; $\times$~=not supported.}
\label{tab:comparison}
\centering
\scriptsize
\begin{tabular}{lcccc}
\toprule
\textbf{System} & \textbf{No Pwd} & \textbf{No PII} & \textbf{Sel.Disc.} & \textbf{Cont.Auth} \\
\midrule
Password + Session & $\times$ & $\times$ & $\times$ & $\times$ \\
OAuth 2.0 (pure)~\cite{oauth2}    & $\approx$ & $\times$ & $\times$ & $\times$ \\
WebAuthn            & \checkmark & $\approx$ & $\times$ & $\times$ \\
\textbf{ZK-Auth (OAuth~2.0 + ZKP)}    & \checkmark & \checkmark & \checkmark & \checkmark \\
\bottomrule
\end{tabular}
\end{table}

The ``OAuth 2.0 (pure)'' row above is the conventional case where OAuth's
Authorization Code grant authenticates the user via a password form;
\textsc{ZK-Auth} retains the same delegation contract
(Section~\ref{sec:oauth_zkp_synopsis}) but replaces that password step
with a Groth16 proof, which is why it is listed separately rather than
as a variant of the pure-OAuth row. Chapter~4 of the full thesis expands
this comparison into a seven-paradigm baseline analysis (adding MFA/TOTP,
PKI smart-card, SAML, and behavioral-biometric-only systems) with
literature-reported quantitative latency figures alongside this
qualitative feature matrix.

%==============================================================================
\chapter{Conclusion and Future Scope}
%==============================================================================

\section{Summary of Findings}

This thesis presented \textsc{ZK-Auth}, a deployable authentication
infrastructure implemented as an \textbf{OAuth~2.0 Authorization Server
in which the conventional password-login step is replaced by a Groth16
zero-knowledge proof}, combining a 932-constraint Groth16 circuit
over BN254, Poseidon Merkle selective disclosure, LSTM behavioral
biometric session continuity trained on authentic public telemetry, and
multi-tenant Ed25519 institutional credential issuance. Empirical
evaluation across four heterogeneous deployment environments demonstrated
a median end-to-end authentication latency of 132.51\,ms---23\% faster
than an Argon2id baseline---with proof generation remaining within
sub-second bounds across the full measured hardware range (61.93\,ms to
361.0\,ms). The behavioral subsystem achieved AUC-ROC~0.8157 with a
6.15\% false-acceptance rate on the public Balabit dataset, and a formal
extraction-based reduction established that the implemented circuit's
shared-witness construction preserves Groth16's soundness guarantee for
its specific Poseidon-bound commitment and nullifier bindings. A
baseline comparison against seven existing authentication paradigms
showed \textsc{ZK-Auth} is competitive with or faster than every
literature-reported figure while uniquely combining password elimination,
PII minimization, selective disclosure, and continuous authentication.

\section{Limitations and Future Scope}

The behavioral model currently trains on mouse-only telemetry; multi-modal
extension incorporating keyboard and scroll signals is identified as
near-term future work. The public Hermez trusted-setup ceremony is used
rather than a deployment-specific one. Future directions include
completing the Merkle-depth and LSTM-window ablations with measured
(rather than derived) data at all settings, iOS WKWebView benchmarking,
FPGA-accelerated proving for edge devices, threshold multi-party signing
for institutional key custody, and integration with a complementary
attribute-based encryption access-control layer.

\section{Concluding Remarks}

These results demonstrate that zero-knowledge-proof-based authentication
is practically deployable at production-acceptable latency across the
realistic spread of consumer client hardware, with direct implications for
cloud security architecture, decentralized identity infrastructure, and
regulatory compliance under India's DPDP Act~2023 and the EU GDPR.

%==============================================================================
% BIBLIOGRAPHY (trimmed to sources directly cited above)
%==============================================================================
\clearpage
\addcontentsline{toc}{chapter}{Bibliography}
\begin{thebibliography}{30}

\bibitem{gmr89}
S. Goldwasser, S. Micali, and C. Rackoff, ``The knowledge complexity of interactive
proof systems,'' \textit{SIAM J. Computing}, vol.~18, no.~1, pp.~186--208, 1989.

\bibitem{goldreich94}
O. Goldreich and Y. Oren, ``Definitions and properties of zero-knowledge proof
systems,'' \textit{J. Cryptology}, vol.~7, no.~1, pp.~1--32, 1994.

\bibitem{bfm88}
M. Blum, P. Feldman, and S. Micali, ``Non-interactive zero-knowledge and its
applications,'' in \textit{Proc. 20th ACM STOC}, 1988, pp.~103--112.

\bibitem{fiatshamir}
A. Fiat and A. Shamir, ``How to prove yourself: Practical solutions to identification
and signature problems,'' in \textit{Proc. CRYPTO 1986}, LNCS vol.~263, 1987,
pp.~186--194.

\bibitem{parno13}
B. Parno, J. Howell, C. Gentry, and M. Raykova, ``Pinocchio: Nearly practical
verifiable computation,'' in \textit{Proc. IEEE S\&P}, 2013, pp.~238--252.

\bibitem{gennaro13}
R. Gennaro, C. Gentry, B. Parno, and M. Raykova, ``Quadratic span programs and
succinct NIZKs without PCPs,'' in \textit{Proc. EUROCRYPT 2013}, LNCS vol.~7881,
pp.~626--645.

\bibitem{groth16}
J. Groth, ``On the size of pairing-based non-interactive arguments,'' in
\textit{Advances in Cryptology -- EUROCRYPT 2016}, LNCS vol.~9666, Springer,
pp.~305--326, 2016.

\bibitem{zcash14}
E. Ben-Sasson et al., ``Zerocash: Decentralized anonymous payments from Bitcoin,''
in \textit{Proc. IEEE S\&P}, 2014, pp.~459--474.

\bibitem{plonk}
A. Gabizon, Z. J. Williamson, and O. Ciobanu, ``PLONK: Permutations over
Lagrange-bases for Oecumenical Noninteractive arguments of Knowledge,''
\textit{IACR ePrint} 2019/953, 2019.

\bibitem{starks}
E. Ben-Sasson et al., ``Scalable, transparent, and post-quantum secure computational
integrity,'' \textit{IACR ePrint} 2018/046, 2018.

\bibitem{dandge2025zkpcloud}
A. Dandge, S. Prasad, N. Tiwari, and M. Chawla, ``Zero-Knowledge Proof for
Authentication in Cloud Computing: A Review,'' in \textit{Proc. FrontSci 2025},
Springer Nature, 2025.

\bibitem{khandait2024}
A. U. Khandait and T. R. Pattanshetti, ``Lightweight authentication system based
on ZKP for IoT devices,'' in \textit{Proc. GCCIT 2024}, IEEE, 2024.

\bibitem{pathak2021}
A. Pathak et al., ``Secure authentication using zero knowledge proof,''
in \textit{Proc. ASIANCON 2021}, IEEE, 2021.

\bibitem{butt2024}
S. A. Butt et al., ``if-ZKP: Intel FPGA-based acceleration of zero knowledge
proofs,'' \textit{arXiv:2412.12481}, 2024.

\bibitem{hou2022}
D. Hou et al., ``Privacy-preserving energy trading using blockchain and ZKP,''
in \textit{Proc. IEEE Blockchain}, 2022.

\bibitem{verma2024}
P. Verma, V. Tripathi, and B. Pant, ``ZeroMedChain: ZKP integration for healthcare
identity and access management,'' in \textit{Proc. INDIACom}, IEEE, 2024.

\bibitem{naidu2023}
D. Naidu et al., ``Smart contract for privacy preserving authentication using ZKP,''
in \textit{Proc. OCIT}, IEEE, 2023.

\bibitem{prasad2025abereview}
S. Prasad, A. Dandge, N. Tiwari, and M. Chawla, ``A Comprehensive Review of
Zero-Knowledge Proofs in Attribute-Based Encryption Frameworks,'' accepted,
\textit{Proc. ICNDIA 2026}, IEEE, 2026.

\bibitem{bbs}
D. Boneh, X. Boyen, and H. Shacham, ``Short group signatures,'' in
\textit{Proc. CRYPTO 2004}, LNCS vol.~3152, Springer, pp.~41--55, 2004.

\bibitem{sdjwt}
O. Terbu, D. Fett, and B. Campbell, ``Selective Disclosure for JWTs (SD-JWT),''
IETF Internet-Draft, 2024.

\bibitem{cl}
J. Camenisch and A. Lysyanskaya, ``An efficient system for non-transferable anonymous
credentials with optional anonymity revocation,'' in \textit{Proc. EUROCRYPT 2001},
pp.~93--118, 2001.

\bibitem{w3cvc}
W3C, ``Verifiable Credentials Data Model v2.0,'' W3C Recommendation, 2024.

\bibitem{w3cdid}
W3C, ``Decentralized Identifiers (DIDs) v1.0,'' W3C Recommendation, 2022.

\bibitem{shepherd95}
S. Shepherd, ``Continuous authentication by analysis of keyboard typing
characteristics,'' in \textit{Proc. ECSD}, 1995, pp.~111--114.

\bibitem{nakkabi11}
A. Nakkabi, I. Trabelsi, and A. Chikh, ``Improving mouse dynamics biometric
performance using variance reduction,'' \textit{IEEE Trans. SMC}, vol.~41, no.~1,
pp.~42--48, 2011.

\bibitem{typenet}
P. Acien et al., ``TypeNet: Deep learning keystroke biometrics,''
\textit{IEEE Trans. BTAS}, vol.~4, no.~1, pp.~57--67, 2022.

\bibitem{typingdna}
TypingDNA, ``Behavioral biometrics authentication,'' Product documentation, 2024.

\bibitem{poseidon}
L. Grassi et al., ``Poseidon: A new hash function for zero-knowledge proof systems,''
in \textit{Proc. USENIX Security}, 2021, pp.~519--535.

\bibitem{bernstein2011}
D. J. Bernstein et al., ``High-speed high-security signatures,'' in \textit{Proc.
CHES 2011}, LNCS vol.~6917, Springer, pp.~124--142, 2011.

\bibitem{hochreiter97}
S. Hochreiter and J. Schmidhuber, ``Long short-term memory,''
\textit{Neural Computation}, vol.~9, no.~8, pp.~1735--1780, 1997.

\bibitem{oauth2}
D. Hardt, ``The OAuth 2.0 Authorization Framework,'' IETF RFC~6749, 2012.

\bibitem{rfc7636}
N. Sakimura, J. Bradley, and N. Agarwal, ``Proof Key for Code Exchange by
OAuth Public Clients,'' IETF RFC~7636, 2015.

\end{thebibliography}

\end{document}